% Problemas del capítulo 3. Generado por tools/problemas/gen.py; no editar a mano.
\section{Problemas}\label{sec:problemas-03}

Graduados: \(\star\) directo, \(\star\star\) para pensar, \(\star\star\star\) desafío. Las respuestas están en el apéndice~\ref{ap:respuestas}; conviene intentarlos antes de mirarlas.

\begin{problema}[\({\star}\)]\label{pb:03-1}
Reflejá \(p=(5,1)\) respecto de la recta \(y=x\) y respecto de la recta \(x+y=2\).
\end{problema}

\begin{problema}[\({\star}{\star}\)]\label{pb:03-2}
Dos espejos pasan por el origen: uno a lo largo del eje \(x\) y otro a \(25^\circ\) de él. Reflejá \((1,0)\) primero en el eje \(x\) y después en el otro espejo. ¿Qué transformación es la composición?
\end{problema}

\begin{problema}[\({\star}{\star}\)]\label{pb:03-3}
Una similaridad \(z\mapsto az+b\) lleva \(0\mapsto 1+i\) y \(2\mapsto 1+5i\). Encontrá \(a\) y \(b\), decí cuánto gira y cuánto agranda, y calculá la imagen de \(1+i\).
\end{problema}

\begin{problema}[\({\star}\)]\label{pb:03-4}
Un plano a escala \(1:500\) se fotocopia ampliado al \(141\)\,\%. ¿Qué escala tiene la copia? ¿Por cuánto se multiplican las áreas dibujadas?
\end{problema}

\begin{problema}[\({\star}{\star}{\star}\)]\label{pb:03-5}
Reflejá en el eje \(x\) y después trasladá \((3,0)\) (una reflexión con deslizamiento). ¿Tiene puntos fijos? ¿Qué da al aplicarla dos veces? Calculá la imagen de \((0,2)\).
\end{problema}
